\documentclass[12pt]{article}
\usepackage{amsmath,amssymb,booktabs,hyperref,geometry}
\geometry{margin=1in}

\newcommand{\vE}{\vec{E}}
\newcommand{\vB}{\vec{B}}

\title{Verification of the Euler--Heisenberg Derivation $\kappa = 3/11$}
\author{DFD Cross-Reference Audit}
\date{23 March 2026}

\begin{document}
\maketitle

%======================================================================
\section{Claim under review}
%======================================================================

The DFD vacuum-loading framework requires a parameter $\kappa$ that
quantifies the asymmetry between the electric ($\epsilon$) and magnetic
($\mu$) response of the vacuum to the scalar field $\psi$.  One
proposed route sets
\begin{equation}
\kappa_{\rm EH} = \frac{7 - 4}{7 + 4} = \frac{3}{11} \approx 0.2727,
\label{eq:claim}
\end{equation}
where the integers 4 and 7 are the coefficients appearing in the
weak-field expansion of the Euler--Heisenberg (EH) effective
Lagrangian of QED.

This document verifies or falsifies each step of that claim.

%======================================================================
\section{Check 1: The exact Euler--Heisenberg coefficients}
%======================================================================

\subsection{Standard form}

The one-loop (spinor-QED) EH effective Lagrangian, expanded to
leading order in $E/E_{\rm cr}$ and $B/B_{\rm cr}$, is
\begin{equation}
\boxed{
\mathcal{L}_{\rm EH}
= \frac{m_e^4}{360\pi^2}
  \left(\frac{e}{m_e^2}\right)^{\!4}
  \bigl[4\,\mathcal{F}^2 + 7\,\mathcal{G}^2\bigr],}
\label{eq:EH}
\end{equation}
where the two Lorentz invariants are
\begin{equation}
\mathcal{F} = \tfrac{1}{2}(\vB^2 - \vE^2),
\qquad
\mathcal{G} = -\vB\cdot\vE.
\end{equation}
Equivalently, using the covariant field-strength tensor,
$\mathcal{F} = \tfrac{1}{4}F_{\mu\nu}F^{\mu\nu}$ and
$\mathcal{G} = \tfrac{1}{4}F_{\mu\nu}\tilde{F}^{\mu\nu}$.

\subsection{Verification of the 4 and 7}

\textbf{Status: CONFIRMED.}

The coefficients 4 and 7 are unambiguously established in the
literature.  The original Heisenberg--Euler (1936) calculation gives
exactly these values for spin-$\tfrac{1}{2}$ QED.  Modern references
confirming this include:
\begin{itemize}
\item Dunne, ``Heisenberg-Euler Effective Lagrangians: Basics and
  Extensions'' (hep-th/0406216), Eq.~(2.19).
\item Dittrich \& Gies, \textit{Probing the Quantum Vacuum} (2000).
\item The ar5iv-rendered version of arXiv:1912.11698, Eq.~(10),
  which explicitly writes
  $\tfrac{m_e^4}{360\pi^2}(e/m_e^2)^4[4\mathcal{F}^2 +
  7\mathcal{G}^2]$.
\item Heinzl et al., Phys.\ Rev.\ D 89, 125003 (2014)
  (arXiv:1312.6419), which identifies $c_+ = 7\alpha/(45\pi E_S^2)$,
  $c_- = 4\alpha/(45\pi E_S^2)$.
\end{itemize}

For scalar QED (spin-0 loop), the coefficients change to $1:1$
(no asymmetry).  The 4:7 asymmetry is a \textbf{fermion-specific}
result, arising from the spin-orbit coupling of the Dirac field
in the loop.

\subsection{Physical origin of the asymmetry}

The asymmetry $4 \neq 7$ has a clear physical origin: the Dirac
spinor couples differently to electric and magnetic backgrounds
because the Dirac equation mixes large and small components
asymmetrically under $\vE$ vs.\ $\vB$.  In the proper-time
(Schwinger) formalism, the eigenvalues of the proper-time
Hamiltonian depend on the two invariants $\mathcal{F}$ and
$\mathcal{G}$ through different secular equations, producing the
4:7 ratio.

A scalar (spin-0) particle in the loop gives equal coefficients,
i.e.\ the 4:7 ratio would become $1:1$, confirming that the
asymmetry is a \emph{spin effect}.

%======================================================================
\section{Check 2: The convention for $\kappa$}
%======================================================================

Two distinct conventions have appeared:

\begin{center}
\begin{tabular}{lcc}
\toprule
\textbf{Convention} & \textbf{Formula} & \textbf{Value} \\
\midrule
``Fractional weight'' & $\kappa = \dfrac{4}{4+7} = \dfrac{4}{11}$ & $0.3636$ \\[8pt]
``Asymmetry parameter'' & $\kappa = \dfrac{7 - 4}{7 + 4} = \dfrac{3}{11}$ & $0.2727$ \\
\bottomrule
\end{tabular}
\end{center}

The DFD vacuum-loading formalism uses the \textbf{asymmetry-parameter}
convention:
\begin{equation}
\kappa \;\equiv\; \frac{c_{\mathcal{G}} - c_{\mathcal{F}}}
                      {c_{\mathcal{G}} + c_{\mathcal{F}}},
\end{equation}
where $c_{\mathcal{F}} = 4$ (the $\mathcal{F}^2$ coefficient) and
$c_{\mathcal{G}} = 7$ (the $\mathcal{G}^2$ coefficient).

\textbf{Rationale for the sign:}  $\kappa > 0$ means the
pseudo-scalar ($\vE\cdot\vB$) sector is more strongly modified than
the scalar ($\vE^2 - \vB^2$) sector.  Equivalently, in the presence
of a pure magnetic background ($\vE = 0$), the coefficient 7
governs the $\vE$-induced perturbation (since $\mathcal{G} =
-\vB\cdot\vE$ mixes $\vE$ into $\vB$), while the coefficient 4
governs the purely magnetic self-interaction.

\textbf{Status:} The formula $\kappa = 3/11$ is \textbf{arithmetically
correct} given the asymmetry-parameter convention and the verified
coefficients 4 and 7.

%======================================================================
\section{Check 3: Energy scale --- does EH apply to DFD?}
%======================================================================

\subsection{EH validity regime}

The Euler--Heisenberg Lagrangian is a low-energy effective theory
valid when the photon energy satisfies
\begin{equation}
\omega \ll m_e c^2 \approx 0.511~\text{MeV},
\end{equation}
and the field strength satisfies $E \ll E_{\rm cr} = m_e^2 c^3 /
(e\hbar) \approx 1.3 \times 10^{18}~\text{V/m}$.

\subsection{DFD vacuum-loading regime}

The DFD scalar field $\psi$ is a \emph{static} or slowly-varying
gravitational-strength background.  The relevant photon energies
for DFD predictions (cavity experiments, atomic clocks) are
\begin{equation}
\omega_{\rm cavity} \sim \text{GHz} \sim 10^{-5}~\text{eV}
\quad \Longrightarrow \quad
\omega / m_e \sim 10^{-11}.
\end{equation}
This is \textbf{deeply within the EH validity regime}.

\subsection{Assessment}

\textbf{Status: CONSISTENT.}  DFD operates at frequencies and field
strengths vastly below the EH cutoff.  The EH effective Lagrangian
is the correct low-energy description of QED vacuum effects in this
regime.  No additional UV physics (heavy particles, string effects,
etc.)\ is needed.

\textbf{Caveat:}  The EH Lagrangian describes the response of the QED
vacuum to \emph{electromagnetic} perturbations.  Whether the $\psi$
field (a gravitational scalar) couples to the vacuum in the same way
as an electromagnetic probe is a separate question addressed in
Check~4.

%======================================================================
\section{Check 4: From photon--photon to the $\epsilon/\mu$ split}
%======================================================================

This is the \textbf{weakest link} in the derivation chain.  Let us
examine it carefully.

\subsection{What the EH Lagrangian directly gives}

The EH Lagrangian describes \emph{photon self-interaction} via a
fermion loop.  In a magnetic background $\vB_0$, it predicts two
refractive indices:
\begin{align}
n_\perp &= 1 + \frac{7\alpha}{90\pi}\,
  \frac{B_0^2}{B_{\rm cr}^2}\sin^2\theta,
\label{eq:nperp}\\
n_\parallel &= 1 + \frac{4\alpha}{90\pi}\,
  \frac{B_0^2}{B_{\rm cr}^2}\sin^2\theta.
\label{eq:npar}
\end{align}
The 7 appears in $n_\perp$ (electric field perpendicular to $\vB_0$)
and the 4 in $n_\parallel$ (electric field in the $\vB_0$--$\vec{k}$
plane).

Since $n^2 = \epsilon_{\rm eff} \cdot \mu_{\rm eff}$, the
birefringence demonstrates that the QED vacuum is an
\emph{anisotropic medium} with distinct effective permittivity and
permeability tensors.

\subsection{What DFD needs}

DFD needs $\kappa$ to describe how a \emph{scalar gravitational
perturbation} $\psi$ modifies $\epsilon_0$ and $\mu_0$ differently:
\begin{equation}
\epsilon(\psi) = \epsilon_0(1 + (1+\kappa)\lambda\psi + \cdots),
\qquad
\mu(\psi) = \mu_0(1 + (1-\kappa)\lambda\psi + \cdots).
\end{equation}
The product $\epsilon\mu$ is constrained by the speed of light, so
the \emph{sum} of the shifts is fixed; $\kappa$ parametrizes how the
shift is \emph{distributed} between the electric and magnetic
sectors.

\subsection{The logical gap}

The EH result tells us the QED vacuum has a built-in
electric/magnetic asymmetry with ratio 7:4.  The \emph{identification}
\begin{equation}
\kappa_{\rm DFD} \;\stackrel{?}{=}\; \kappa_{\rm EH}
\label{eq:identification}
\end{equation}
requires an additional assumption:

\begin{quote}
\textbf{Assumption:} The scalar field $\psi$ couples to the
electromagnetic vacuum through the same fermion-loop polarization
tensor that generates the EH effective Lagrangian.  The coupling
$\mathcal{L}_{\rm int} = -\tfrac{1}{2}\lambda\psi T^\mu{}_\mu$
projects onto the electric and magnetic sectors with the same
relative weights as the EH coefficients.
\end{quote}

This assumption is \textbf{plausible but not proven}.  The trace
anomaly gives
$\langle T^\mu{}_\mu\rangle = (\alpha/6\pi)(\vB^2 - \vE^2)$,
which shows the $\vB$ sector dominates (consistent with $\kappa > 0$
and with the 7 > 4 pattern), but does not uniquely fix $\kappa =
3/11$.  The trace anomaly is a \emph{single} number; extracting the
\emph{ratio} of two components requires knowing the full tensor
decomposition of the $\psi$--vacuum coupling.

\subsection{Assessment}

\textbf{Status: PLAUSIBLE BUT NOT RIGOROUS.}

The EH coefficients correctly describe the electromagnetic
vacuum's intrinsic electric/magnetic asymmetry.  The identification
$\kappa_{\rm DFD} = 3/11$ is a \emph{natural} choice (it is the
unique dimensionless ratio built from 4 and 7), but it assumes that
the $\psi$--vacuum coupling inherits the full EH structure.

An alternative (and arguably cleaner) route through the electroweak
sector gives $\kappa = \cos 2\theta_W \approx 0.538$, which differs
by a factor $\sim 2$.

%======================================================================
\section{Check 5: Assumptions beyond QED + vacuum loading}
%======================================================================

The derivation $\kappa = 3/11$ requires:

\begin{enumerate}
\item \textbf{Standard QED at one loop:} The EH Lagrangian with
  coefficients 4 and 7.  (\emph{Verified.})

\item \textbf{Fermion dominance:} Only spin-$\tfrac{1}{2}$ particles
  contribute.  Scalar loops give 1:1 (no asymmetry); vector boson
  loops give different ratios.  If multiple species contribute, the
  effective $\kappa$ is a weighted average.  (\emph{Reasonable for
  pure QED; less clear in the full Standard Model.})

\item \textbf{Weak-field limit:} The expansion to order $(\alpha
  E/E_{\rm cr})^2$ is used.  Higher-order terms would modify
  $\kappa$.  (\emph{Justified for any laboratory or astrophysical
  field.})

\item \textbf{Scalar coupling identification:} The $\psi$ field
  couples to the vacuum polarization tensor with the same relative
  weights as the EH effective Lagrangian.
  (\emph{\textbf{Assumed, not derived.}  This is the key additional
  input beyond QED + vacuum loading.})

\item \textbf{No electroweak corrections:} The derivation uses pure
  QED.  Including $W$/$Z$ loops would modify the effective
  coefficients.
\end{enumerate}

%======================================================================
\section{Comparison with the electroweak route}
%======================================================================

The DFD framework also admits
\begin{equation}
\kappa_{\rm EW} = \cos 2\theta_W = 1 - 2\sin^2\theta_W
\approx 0.538.
\end{equation}
This is derived from the electroweak mixing angle, which controls
how the photon inherits its coupling from the $B$ (hypercharge)
and $W^3$ (isospin) gauge bosons after symmetry breaking.

\begin{center}
\begin{tabular}{lccl}
\toprule
\textbf{Route} & $\kappa$ & \textbf{Numerical} & \textbf{Status} \\
\midrule
Euler--Heisenberg & $3/11$ & 0.2727 & Coefficients verified;
  coupling identification assumed \\
Electroweak mixing & $\cos 2\theta_W$ & 0.5376 & Clean derivation;
  probes EW vacuum \\
Trace anomaly & --- & $> 0$ & Sign only; no unique value \\
\bottomrule
\end{tabular}
\end{center}

The two routes disagree by a factor $\sim 2$ because they probe
\emph{different physics}:
\begin{itemize}
\item The EH route probes the \emph{QED vacuum} (electron loop only).
\item The EW route probes the \emph{electroweak vacuum} (full gauge
  sector after symmetry breaking).
\end{itemize}
Which is correct depends on whether $\psi$ couples to the pure QED
vacuum or to the full electroweak vacuum---a question that experiment
can resolve via the predicted TE/TM cavity splitting.

%======================================================================
\section{Verdict}
%======================================================================

\begin{enumerate}
\item \textbf{Coefficients 4 and 7:} \textsc{Verified}.  These are
  the standard, well-established one-loop spinor-QED
  Euler--Heisenberg coefficients.

\item \textbf{Arithmetic $\kappa = (7-4)/(7+4) = 3/11$:}
  \textsc{Verified}.  This is correct in the asymmetry-parameter
  convention.

\item \textbf{Energy-scale consistency:}  \textsc{Verified}.  DFD
  operates deep within the EH validity regime ($\omega/m_e \sim
  10^{-11}$).

\item \textbf{Connection to $\epsilon/\mu$ split:}
  \textsc{Plausible but not rigorous}.  The identification of the EH
  asymmetry with the DFD vacuum-loading split requires an assumption
  about how $\psi$ couples to the QED vacuum polarization tensor.
  This assumption is physically motivated but not derived from first
  principles within DFD.

\item \textbf{Assumptions beyond QED + vacuum loading:}  The key
  extra assumption is the scalar-coupling identification (item 4 in
  Check~5).  Additionally, the derivation assumes pure QED (no
  electroweak corrections), which may not be the full story.
\end{enumerate}

\subsection*{Bottom line}

The derivation $\kappa = 3/11 \approx 0.27$ is
\textbf{mathematically correct} and \textbf{physically well-motivated},
but it is \textbf{not the unique prediction} of the DFD framework.
The electroweak route gives $\kappa \approx 0.54$.  The two values
are experimentally distinguishable via TE/TM cavity splitting or
LPI violation measurements.  \textbf{The EH coefficients 4 and 7 are
solid; the question is whether the EH photon-photon asymmetry maps
onto the DFD $\psi$--vacuum constitutive split.}

\end{document}
